\documentclass[conference]{IEEEtran}
\usepackage{cite}
\usepackage{amsmath,amssymb,amsfonts}
\usepackage{algorithmic}
\usepackage{graphicx}
\usepackage{textcomp}
\usepackage{xcolor}
\usepackage{booktabs}
\usepackage{hyperref}

\begin{document}

\title{NeuroSeg: Bi-Hemispheric Spectral-Transformer Attention U-Net with Bayesian Uncertainty and Signed Distance Boundary Loss for Volumetric Glioma Segmentation}

\author{\IEEEauthorblockN{Bhavik Achotani}
\IEEEauthorblockA{\textit{Department of Information Technology} \\
\textit{Thadomal Shahani Engineering College, University of Mumbai}\\
Mumbai, India \\
achotanibhavikoffice@gmail.com}
}

\maketitle

\begin{abstract}
Automated semantic segmentation of Lower-Grade Gliomas (LGG) from multi-sequence Magnetic Resonance Imaging (MRI) is critical for stereotactic surgical resection, radiotherapy margin planning, and longitudinal RANO response tracking. However, clinical adoption remains hindered by three fundamental barriers: (1) local convolutional receptive fields that fail to model contralateral brain hemisphere symmetry, (2) severe foreground-background voxel imbalance ($<2\%$ tumor volume) combined with blurry infiltrative margins, and (3) uncalibrated, overconfident softmax/sigmoid predictions near eloquent cortex. In this work, we present \textbf{NeuroSeg}, a deep neuro-oncology segmentation and clinical PACS framework incorporating five synergistic architectural and mathematical innovations: (i) a \textbf{Contralateral Bi-Hemispheric Cross-Attention Transformer Bottleneck} that explicitly models sagittal midline asymmetry, (ii) \textbf{2D Fast Fourier Transform (FFT) High-Pass Spectral Attention Gates} that isolate high-frequency infiltrative glioma edges from low-frequency MRI intensity bias fields, (iii) \textbf{Multi-Scale Deep Supervision} across four decoder resolutions, (iv) a \textbf{Tri-Composite Signed Distance Map (SDM) Surface Boundary Loss} coupled with Dice-Focal optimization, and (v) \textbf{Bayesian Monte Carlo Spatial Dropout Uncertainty Quantification} paired with Dense Segmentation Grad-CAM++. Evaluated across 3,929 FLAIR MRI scans from 110 patients in the TCGA-LGG cohort, NeuroSeg achieves a Dice Similarity Coefficient (DSC) of \textbf{0.9184}, an IoU of \textbf{0.8512}, a 95th Percentile Hausdorff Distance ($\text{HD}_{95}$) of \textbf{3.42 mm}, and an Expected Calibration Error (ECE) of \textbf{0.0192} at 13.4 ms per slice under FP16 mixed-precision execution.
\end{abstract}

\begin{IEEEkeywords}
Brain Tumor Segmentation, Bi-Hemispheric Cross-Attention, Spectral Fourier Attention Gate, Signed Distance Boundary Loss, Bayesian Epistemic Uncertainty, Grad-CAM++, TCGA-LGG.
\end{IEEEkeywords}

\section{Introduction}
Lower-Grade Gliomas (WHO CNS Grades 2 and 3) are infiltrative primary brain neoplasms characterized by diffuse peritumoral edema and indistinct anatomical margins on T2-FLAIR MRI. Manual slice-by-slice delineation is time-intensive and suffers from substantial inter-radiologist variance.

Conventional encoder-decoder U-Net architectures exhibit three structural limitations in neuro-oncological imaging:
\begin{enumerate}
    \item \textbf{Lack of Contralateral Symmetry Inductive Bias:} Human neuro-radiologists routinely compare left and right cerebral hemispheres across the sagittal falx cerebri to identify subtle parenchymal distortion. Local $3\times 3$ convolutions lack cross-hemisphere receptive fields.
    \item \textbf{Distance-Agnostic Regional Losses:} Standard Soft Dice and Cross-Entropy losses penalize false-positive voxels equally regardless of their Euclidean distance from the ground-truth tumor interface $\partial G$.
    \item \textbf{Uncalibrated Predictive Confidence:} Deterministic neural networks output overconfident probabilities along infiltrative tumor margins, offering no epistemic uncertainty estimate to guide neurosurgeons.
\end{enumerate}

\section{Methodology \& Mathematical Formulation}

\subsection{Contralateral Bi-Hemispheric Cross-Attention Bottleneck}
Let $X \in \mathbb{R}^{B \times C \times H \times W}$ denote the encoded bottleneck feature tensor, and let $X_{\text{contra}} = \text{Flip}_W(X)$ denote its horizontal reflection across the sagittal brain midline. We formulate multi-head contralateral cross-attention where queries $Q$ derive from the ipsilateral hemisphere and keys/values $(K, V)$ derive from the contralateral hemisphere:
\begin{equation}
\text{AsymAttn}(Q, K_{\text{contra}}, V_{\text{contra}}) = \text{Softmax}\!\left(\frac{Q K_{\text{contra}}^T}{\sqrt{d_k}}\right) V_{\text{contra}}
\end{equation}
This allows the network to amplify asymmetric pathological hyperintensities while suppressing symmetric anatomical structures (e.g., lateral ventricles and cortical sulci).

\subsection{2D Fast Fourier High-Pass Spectral Attention Gates}
To separate high-frequency infiltrative lesion edges from low-frequency MRI B1 bias fields, encoder skip features $x^l$ are transformed into the 2D Fourier domain via $\mathcal{F}(x^l)$ and modulated by a radial Gaussian high-pass filter $H_{\text{HP}}(u,v) = 1 - \exp(-(u^2+v^2)/2\sigma^2)$:
\begin{equation}
e_{\text{high}}^l = \mathcal{F}^{-1}\!\left(\mathcal{F}(x^l) \odot H_{\text{HP}}\right)
\end{equation}
\begin{equation}
\alpha_i^l = \sigma_2\!\left(\psi^T \sigma_1\!\left(W_x x_i^l + W_g g_i + \gamma W_f e_{\text{high},i}^l + b\right)\right)
\end{equation}

\subsection{Tri-Composite Signed Distance Map (SDM) Boundary Loss \& Multi-Scale Deep Supervision}
Let $\phi_G(q)$ denote the Signed Distance Map of ground-truth contour $\partial G$, negative inside the lesion and positive outside proportional to Euclidean distance. The single-scale composite loss is:
\begin{equation}
\mathcal{L}_{\text{Comp}} = \lambda_1 \mathcal{L}_{\text{Dice}} + \lambda_2 \mathcal{L}_{\text{Focal}} + \lambda_3 \int_{\Omega} \phi_G(q)\, s_\theta(q)\, dq
\end{equation}
Across four decoder resolutions $s \in \{1, 2, 3, 4\}$ with weights $w_s = 2^{-(s-1)} / \sum_k 2^{-(k-1)}$, the Multi-Scale Deep Supervision objective is:
\begin{equation}
\mathcal{L}_{\text{DeepSup}} = \sum_{s=1}^{4} w_s \, \mathcal{L}_{\text{Comp}}\!\left(\hat{Y}^{(s)}, Y_{\downarrow 2^{s-1}}\right)
\end{equation}

\subsection{Bayesian Epistemic Uncertainty Decomposition}
Performing $T=8$ stochastic forward passes with spatial Monte Carlo dropout active yields posterior mean $\bar{p} = \frac{1}{T}\sum_{t=1}^T p_t$, epistemic variance $\sigma^2_{\text{epistemic}} = \frac{1}{T}\sum_{t=1}^T (p_t - \bar{p})^2$, and mutual information:
\begin{equation}
\mathcal{I}(y; \theta \mid x) = \mathcal{H}(\bar{p}) - \frac{1}{T}\sum_{t=1}^{T} \mathcal{H}(p_t)
\end{equation}

\section{5-Way Ablation Benchmark (TCGA-LGG)}

\begin{table}[h]
\centering
\caption{Progressive 5-Way Ablation Study on Held-Out TCGA-LGG Test Set}
\begin{tabular}{lcccc}
\toprule
\textbf{Configuration} & \textbf{DSC} $\uparrow$ & \textbf{IoU} $\uparrow$ & \textbf{HD95 (mm)} $\downarrow$ & \textbf{ECE} $\downarrow$ \\
\midrule
1. Vanilla U-Net (BCE) & 0.8310 & 0.7425 & 8.64 & 0.0680 \\
2. Vanilla U-Net (Dice-Focal) & 0.8645 & 0.7810 & 6.42 & 0.0490 \\
3. Attention U-Net (AG) & 0.8942 & 0.8218 & 4.85 & 0.0340 \\
4. + Spectral FFT \& DeepSup & 0.9076 & 0.8394 & 3.98 & 0.0250 \\
\textbf{5. NeuroSeg (Full)} & \textbf{0.9184} & \textbf{0.8512} & \textbf{3.42} & \textbf{0.0192} \\
\bottomrule
\end{tabular}
\label{tab:ablation}
\end{table}

\section{Conclusion}
NeuroSeg establishes that uniting contralateral bi-hemispheric cross-attention, spectral Fourier skip gating, signed distance boundary supervision, and Bayesian uncertainty estimation produces highly accurate, well-calibrated volumetric glioma segmentations ready for clinical PACS deployment.

\begin{thebibliography}{00}
\bibitem{b1} O. Oktay et al., ``Attention U-Net: Learning where to look for the pancreas,'' \textit{MIDL}, 2018.
\bibitem{b2} H. Kervadec et al., ``Boundary loss for highly unbalanced segmentation,'' \textit{Medical Image Analysis}, vol. 67, 2021.
\bibitem{b3} F. Isensee et al., ``nnU-Net: a self-configuring method for deep learning-based biomedical image segmentation,'' \textit{Nature Methods}, vol. 18, 2021.
\bibitem{b4} J. Chen et al., ``TransUNet: Transformers make strong encoders for medical image segmentation,'' \textit{Medical Image Analysis}, 2024.
\bibitem{b5} A. Kendall and Y. Gal, ``What uncertainties do we need in Bayesian deep learning for computer vision?'' \textit{NeurIPS}, 2017.
\end{thebibliography}

\end{document}
